% 393_Rauschgrenze_Lernende_Systeme_En_ch.tex
% Johann Pascher, 13 Oct. 2026
% The noise limit and learning systems: questions and answers

\providecommand{\BM}{\textbf{[B]}}
\providecommand{\KM}{\textbf{[K]}}
\providecommand{\SM}{\textbf{[S]}}
\providecommand{\XM}{\textbf{[X]}}
\providecommand{\QM}{\textbf{[Q]}}

\chapter*{The Noise Limit and Learning Systems}
\addcontentsline{toc}{chapter}{The Noise Limit and Learning Systems}

\begin{abstract}
The corpus knows a noise limit: beyond a certain measurement tolerance, harmonic relations can no
longer be distinguished from chance (Doc.~343, Theorem G$'$). In questions and answers, this document
examines whether neural networks and language models run into the same limit, why they nevertheless
almost always produce readable sentences, why language forgives errors and numbers do not, and what
role probability and the weights of a network play. A null test with an exhaustive symbolic search
over about 23 million expressions shows that the dimensionless measured quantities of the basic form
carry no harmonic structure above chance level on their own, and that even a genuine rule such as the
Balmer series becomes visible from a few examples only with the appropriate, narrow prior. The common
core of all answers: discrete structure is error-tolerant, the noise limit sits at the transition from
the discrete to the continuum, and structure that the data do not yield must be specified and then
tested.
\end{abstract}

% ============================================================
\section*{Status markers}
% ============================================================
\begin{center}
\begin{tabular}{@{}lp{11cm}@{}}
\textbf{[B]} & Proved algebraically or information-theoretically. \\[2pt]
\textbf{[K]} & Confirmed numerically (verification script). \\[2pt]
\textbf{[S]} & Setting, open, or classification without a calculation of its own. \\[2pt]
\textbf{[X]} & Excluded by the calculation. \\[2pt]
\textbf{[Q]} & Taken from the literature, with source. \\
\end{tabular}
\end{center}

\noindent The document is structured as a sequence of questions and answers. The questions are kept
in the order in which they arise from one another; each answer stands on its own. Section~14
summarises the common core in a single picture.

% ============================================================
\section{The noise limit in the corpus}
% ============================================================

\textbf{Question.} What is meant by the noise limit?

\textbf{Answer.} A measured ratio is known only up to a tolerance $\varepsilon$. The larger the
admitted denominator, the more densely the fractions lie; beyond the Farey height
\[
  Q_c=\frac{\pi}{\sqrt{6\,\varepsilon}}
\]
the tolerance window contains so many fractions that a genuine ratio can no longer be separated from a
chance one (Doc.~343, Theorem G$'$). For $\varepsilon=1.05\,\%$, the size of the fractal correction,
$Q_c\approx12.5$; for $\varepsilon=100\,\xi\approx1.33\,\%$, $Q_c\approx11.1$ (Doc.~359, Theorem~B)
\KM. Harmonic relations with small denominators remain visible; finer ones are lost in the noise. The
limit is a property of the data, not of the method: what is not contained in the signal cannot be
recovered by any analysis, neither by a Fourier analysis nor by an autocorrelation \BM.

% ============================================================
\section{Do neural networks and language models run into the same limit?}
% ============================================================

\textbf{Question.} Should neural networks and language models not also run into this limit in pattern
recognition?

\textbf{Answer.} Yes. A network is a method like any other; it can reach the limit but not exceed it.
Beyond the limit it no longer recognises but adapts to the noise and takes chance proximity for
structure; in learning theory this is called overfitting. In addition, such networks have a known
property: they learn coarse, low-frequency structures first and fine ones late
\QM\,\cite{Rahaman2019}. Fine structures, however, are precisely those closest to the noise.

Language models are in a less favourable position. They learn from text, not from measured signals,
and have no measurement resolution of their own. With numerical ratios they tend towards false hits:
given enough combinations, a \enquote{harmonic} explanation can almost always be found (look-elsewhere
effect). This is why verification scripts in the corpus are set up as null tests: if a random
selection yields the same result, it is not counted. This applies to every suggestion made by a
language model.

Networks become better than the raw data only through built-in prior knowledge. Doc.~359 Theorem~D
shows this with an example: whether the admissibility criterion \enquote{prime factors only from
$\{2,3,5,7,11,13\}$} is learnable is decided by the hypothesis class. Without a structural assumption
the learning bound requires about 84\,000 labelled examples; with the structural assumption
\enquote{prime factors from a fixed set} at most 407, and from positive examples alone the correct
prime set is then reliably recovered after only 30 examples \KM\,\cite{ShalevShwartz2014}.

% ============================================================
\section{Why do language models still produce readable sentences?}
% ============================================================

\textbf{Question.} If this is so, why do language models hardly ever produce unreadable sentences?

\textbf{Answer.} Because grammar and sentence structure are what occurs most densely in the training
data. There are trillions of words, the same rules recur in millions of different sentences, there is
no measurement noise, and the architecture is tailored to sequences of words. This is exactly the
favourable case of Doc.~359 Theorem~D: a suitable prior and plenty of examples.

The models become weak where the pattern is rare or exact: precise numbers, chains of calculation,
rare facts, bibliographic references. There they produce something plausible, and the sentence
remains fluent. Readability is the dense pattern, correctness often the sparse one. A well-written
sentence from a language model therefore says nothing about the correctness of its content.

% ============================================================
\section{Should harmonic patterns not be as dense as grammar?}
% ============================================================

\textbf{Question.} If FFGFT is correct, harmonic ratios pervade all of nature. Should they then not
occur just as densely as grammar?

\textbf{Answer.} Dense in occurrence, yes; dense in variety, no. For learning a rule, what counts is
the number of \emph{different} examples, not the number of repetitions. The electron occurs an
unimaginable number of times, but each time it yields the same mass ratio. The basic structure of
FFGFT lies in a few numbers: some mass ratios, couplings and mixing angles, a few dozen altogether.
From so few values no method can separate a rule from chance.

Where harmonic patterns do occur in great variety, they have also been found: in atomic spectra with
thousands of lines following one integer construction principle, in crystal lattices, in orbital
resonances and in musical instruments. Balmer found the rule of the hydrogen lines in 1885 by hand
from four lines \QM\,\cite{Balmer1885}; symbolic regression has rediscovered known formulae of physics
from data \QM\,\cite{Udrescu2020}.

A second reason comes from the corpus itself. According to FFGFT the measured ratios are not exact
fractions: the basic form is exact, the measured value carries the embedding error (Doc.~333,
Doc.~358 Theorem~A), and the PDG values contain model-dependent contributions (Doc.~351). A learning
procedure that sees only measured values learns the shifted form, not the rule beneath it. In
language this layer does not exist.

% ============================================================
\section{The null test}
\label{sec:nulltest}
% ============================================================

\textbf{Question.} Can it be tested how much structure the measured values carry themselves and how
much the prior contributes?

\textbf{Answer.} Yes, with a symbolic regression in its exact form, the exhaustive search over an
expression grammar. All expressions
\[
  x=c\cdot d^{1/k}\cdot\xi^{m/n}\cdot\pi^{s}
\]
are searched, with a coprime fraction $c=p/q$, $p,q\le200$, roots $d^{1/k}$ with $d\le30$ and
$k\in\{2,3\}$, $\xi=1/7500$ with $n\in\{1,2,3,6\}$ and $|m/n|\le4.5$, and $s\in\{-1,0,1\}$. The prime
factors of $p$, $q$ and $d$ are taken from a chosen set of primes. With $\{2,3,5,7,11,13\}$ these are
23\,084\,325 expressions; separately, the search runs over the fractions $p/q$ alone (3109
expressions). Each expression is assigned a complexity in bits that grows with numerator,
denominator, root and exponent. The $p$-value is the fraction of random numbers (log-uniform, one
decade around the measured value) that the grammar hits equally well at equal or lower complexity.
Hits more precise than the measurement uncertainty count only up to that uncertainty.

\textbf{With prior.} The formulae of the basic form from Doc.~387 are fixed; they are compared with
all expressions of equal complexity in the same grammar:

\begin{center}
\small
\begin{tabular}{@{}lllr@{}}
\toprule
\textbf{Quantity} & \textbf{Formula} & \textbf{Deviation} & \textbf{$p$} \\
\midrule
$v/E_P$ & $\xi^4/(5\pi)$ & $-0.23\,\%$ & $0.02$ \\
$M_W/M_Z$ & $\sqrt{7/9}$ & $+0.06\,\%$ & $0.09$ \\
$\lambda_{\mathrm{CKM}}$ & $\xi^{1/6}$ & $+0.45\,\%$ & $0.11$ \\
$m_h/M_Z$ & $11/8$ & $+0.15\,\%$ & $0.42$ \\
$m_\mu/m_e$, $m_\tau/m_\mu$, $\Delta m^2$ ratio, $\sin^2\theta_{13}$, $m_t/M_Z$ &  & $0.1$--$1.5\,\%$ & $1.0$ \\
\bottomrule
\end{tabular}
\end{center}

\noindent Combined by Fisher's method this gives $p=0.41$ \KM. At deviations of $0.5$ to $1\,\%$ the
grammar hits every number; there the lepton formulae carry no numerical evidential weight. The most
striking case is $v/E_P$. For $\alpha^{-1}=3700/27$ (Doc.~338): the deviation of $7.6$\,ppm is also
reached by a fraction with denominator up to 27 for a random number in about half of the cases
($p\approx0.47$) \KM. This holds, however, only if $3700/27$ is treated as a freely searched
fraction. In the corpus it arises as the reduction of $\tfrac{74}{75}\cdot7500/54$ from $1/\xi$, the
rule $74/75$ and the integer $54$; calculated according to these constituents, the chance
probability is about $10^{-3}$ (Docs.~394, 395) \KM.

\textbf{Without prior.} The blind search over the same nine quantities stays at chance level for every
set of primes: for $\{2,3,5,7,11,13\}$, for the control sets $\{2,3,5,7,17,19\}$ and
$\{2,3,5,13,17,23\}$, for $\{2,3,5,7\}$, and likewise for pure random numbers \KM.

\textbf{Positive control.} The Balmer lines of hydrogen have the exact ratios $27/20$, $28/25$ and
$200/189$; the air wavelengths match them to $10^{-5}$ \QM\,\cite{NIST}. With the narrow grammar of
fractions this rule is found from just three ratios (combined $p=0.002$), with the full grammar it is
not ($p=0.47$) \KM. The 23 million expressions lie so densely that random numbers are hit equally well.

\textbf{What follows.} Even a true rule is visible from a few examples only with the right, narrow
prior. Numerical proximity does not carry FFGFT; that must be done by the chain of derivation:
geometry, Galois structure, formula. This corresponds to the reading of Doc.~384 and Doc.~387,
according to which the ratios of the basic form test the structure. Nothing is refuted by the test;
it measures only what the numbers yield on their own.

\textbf{Limits of the test.} The bit measure of complexity is a choice; a different measure shifts the
individual $p$-values \SM. There are only nine quantities. The formulae were found with knowledge of
the measured values; the comparison with equally complex alternatives is therefore the fair standard,
not the raw hit. The lepton residuals of $0.5$ to $1\,\%$ are regarded in the corpus as genuine
structure ($K_{\mathrm{frak}}$, Doc.~352); the test can only count them as misses.

% ============================================================
\section{Language changes constantly, and yet it works}
% ============================================================

\textbf{Question.} Language changes all the time, new word formations appear, grammar varies across
dialects. Why does it work nevertheless?

\textbf{Answer.} The very changeability confirms the picture. Language operates with a large
tolerance and at the same time provides very many examples.

On tolerance: a sentence remains intelligible if a word is misspelt, a dialect word is substituted or
the word order is twisted. Printed language is highly redundant \QM\,\cite{Shannon1951}. The precision
that language requires is coarse; in the terms of the corpus: large $\varepsilon$, hence small $Q_c$.
Language lives in the range in which structures are still separable.

On the examples: new word formations and dialects are not isolated cases but patterns in their own
right, with many instances. A dialect has millions of speakers and countless sentences; new words
follow old construction rules, so \enquote{to google} is conjugated like any regular verb. A learning
procedure therefore does not learn a single rigid grammar but a distribution of variants, and each
variant is densely attested. Beneath the dialects lies a common deeper structure, as with Balmer: many
lines, one rule. Where it fails, the picture fits as well: with a small dialect that has little text,
a rare technical term or an entirely new expression, language models become uncertain.

Music shows both sides. The harmonic grammar of octave, fifth and third, that is, of the primes 2, 3
and 5, is universally recognisable and has many \enquote{dialects}: just intonation, equal
temperament, meantone, folk-music tunings. It works because the ear tolerates deviations of the size
of a comma. This is the statement of Doc.~342 and Doc.~343: up to $k\approx3$ the Galois structures are
separable at this tolerance, beyond it they are not \BM. The harmonics of the small primes are as dense
and robust in nature as grammar. The fine ratios of particle masses, where a tenth of a per cent
matters, lie beyond this limit; there the structure must come from the derivation.

% ============================================================
\section{Why does language forgive errors and numbers do not?}
% ============================================================

\textbf{Question.} Why does language allow this tolerance, while mathematics and numbers do not?

\textbf{Answer.} Because of redundancy. In terms of communications engineering, language is an
error-correcting code and a number an uncoded transmission.

Most random sequences of letters are not words. If one letter is wrong, the result is usually an
invalid word, and the receiver corrects it to the nearest valid one; syntax and context provide
additional checks. This corresponds to a code with a large Hamming distance: errors are detected and
corrected. Many forms moreover lead to the same meaning, and at the destination only the meaning
counts.

For a number the distance is zero. Every sequence of digits is a valid number; with one wrong digit
$137.036$ becomes $137.037$, not a faulty number but a different, equally valid one. Nothing signals
the error, nothing corrects it. Mathematics also chains: an error enters every following step and is
often amplified, whereas errors in a sentence do not add up.

In counting, by contrast, mathematics is tolerant. Integers, group orders and winding numbers behave
like words: between 3 and 4 there is no valid group order, and $3.01$ is immediately recognised as
invalid. In the corpus this appears as \enquote{state versus process}: counting is exact, a bit is a
winding number (Doc.~343, Doc.~247). The algebraically proved statements of FFGFT, namely
$\mathbb Z_3$, $GF(27)$, orbits and fixed points, are therefore as error-tolerant as language. The
intolerance arises at a single point: at the transition from the discrete to the continuum, when an
integer ratio is compared with a measured value on the real axis. There every neighbouring value is
valid, and no distance absorbs the noise.

% ============================================================
\section{Do AI models draw on a finished framework?}
% ============================================================

\textbf{Question.} AI models ultimately have a finished framework from which they draw; not everything
has to be searched anew for patterns. How does this fit with the noise limit?

\textbf{Answer.} After training a model is indeed a finished framework; when answering, nothing is
searched any more, what has been learned is applied. The weights are the frozen prior, in the sense of
Doc.~359 Theorem~D the hypothesis class.

This framework, however, has itself been learned from data. It is a condensed form of the training
texts and contains only what occurred densely there. Where the texts were sparse it has gaps, and the
model fills them with something plausible; this is where invented references come from. It is
moreover statistical rather than exact, built for language and unreliable for exact numbers. Modern
systems therefore call a tool for the exact part: the statistical framework decides what is to be
calculated, the tool calculates. In the work on the corpus this is the division of labour between
formulation and verification script.

FFGFT is also a finished framework: $\xi$, $T^4/\mathbb Z_3$, the Galois structure. It is, however,
not learned from data but set and derived. The learned framework is broad and tolerant but without a
guarantee of exactness; the set one is exact and discrete but must prove itself against the measured
value, precisely where the discrete meets the continuum. After the null test, a learning procedure
cannot find the FFGFT structure on its own; it can, however, be given it as a built-in framework, as
proposed by the architecture sketch of Doc.~359 \SM.

% ============================================================
\section{Where do the rules of an AI assistant come from?}
% ============================================================

\textbf{Question.} An AI assistant works with many prescribed rules. Are they learned or entered
directly?

\textbf{Answer.} Both, in three layers.

The first layer is learned and makes up the largest part. First, training on very large amounts of
text produces the statistical framework. Then a second training stage shapes the behaviour, with
ratings by humans and with written principles such as honesty, avoidance of harm and the admission of
uncertainty. These principles, too, then reside in the weights and are not looked up anew for every
answer.

The second layer is supplied afresh with every conversation: instructions from the product, for
example how tools are to be used and what may happen only after asking, as well as the user's stored
preferences. This layer is the easiest to change; it takes effect immediately, without new training.
The third layer is the running context: everything that stands in the conversation and everything
that has been read.

The layers can contradict one another. During the work on this corpus, an automatic check demanded
after every answer that changes be transferred to the public repository, while the author's stored
preference forbids exactly that; the explicit preference was followed. The division fits the topic
of this document: the first layer is the learned, tolerant framework, the second and third are set
rules, and the verification scripts are the exact tool that checks both.

% ============================================================
\section{How does the next step arise after an error message?}
% ============================================================

\textbf{Question.} Between an error message from a script and the decision about what is most sensible
next, there must be many layers.

\textbf{Answer.} Yes, but differently from what one might expect: there is no separate decision
module, only layers in the network and a simple loop around it.

The loop: the assistant writes a tool call. The surrounding program executes it and appends the
result, including an error message, to the text of the conversation. Then the whole is read again and
the next step is written. A plan exists only as text in this context.

The layers: the network itself has dozens to more than a hundred consecutive layers, each of which
further processes the entire preceding text. From error message, goal, files read and previous
attempts the next step emerges. That it is sensible is trained: in the second training stage such
sequences were practised on many tasks, and sequences that led to the goal were reinforced.

Two examples from the resetting of the book series: an \enquote{Improper alphabetic constant} message
when typesetting a chapter title containing $\varphi$ is a densely attested pattern; the cause,
mathematics in the PDF bookmarks, is recognised at once. A volume that produced only 18 instead of
more than 300 pages without any error message required an inference instead: the input ended too
early, presumably because of an \verb|\end{document}| in one chapter, which a search confirmed. Where
patterns are sparse, things are tried out and several attempts are needed. The same therefore holds
for the assistant as for FFGFT: try, check exactly, measure with verification scripts, rather than
trusting one's own impression.

% ============================================================
\section{Chance and probability}
% ============================================================

\textbf{Question.} Does the noise limit not also play a role in chance and probabilities?

\textbf{Answer.} Yes, quite directly: probability theory is the mathematics of the noise limit. Two
kinds of probability must be distinguished.

Probability as ignorance: a measured value fluctuates. Averaging $N$ independent measurements reduces
the noise as $1/\sqrt N$; a hundred times as many measurements reduce it to a tenth \BM. This is why
$m_\mu/m_e$ is known to $10^{-8}$, although each individual measurement is coarser. This holds only
for random errors; systematic contributions such as the model assumptions of the PDG values
(Doc.~351) do not average out. The $p$-value of the null test belongs here as well: it is the noise
limit expressed as a probability.

In language models probability is the output itself. For each word a distribution over all possible
next words is computed and one is selected. Where the training data were dense the distribution is
sharp and the answer stable; where they were sparse it is flat, and the selection is almost a matter
of chance. A flat distribution is the noise limit of the model, and from outside it cannot be seen,
because the sentence still sounds fluent.

Probability as structure: in quantum mechanics probabilities are exact numbers. According to Doc.~367
the Koide phase $\theta$ and the icosahedral transition probability $p_0$ are the same quotient
$2/3^2=2/9$, two readings of one Born-rule structure. Such a probability is as discrete as a group
order. Noise enters only with the measurement: a frequency from $N$ events scatters around $2/9$ by
about $\sqrt{p(1-p)/N}$; to separate $2/9$ from the neighbouring fraction $1/5$ at three standard
deviations requires about 3200 events \KM. In this sense Doc.~230 also treats the Bell correlations:
real, but geometrically structured. Probability thus appears twice: as an object of the theory, exact,
and as a tool for judging whether the exact can still be recognised in the noise.

% ============================================================
\section{Does the noise limit enter the weights?}
% ============================================================

\textbf{Question.} What about the weights of a network: do the noise limits enter there?

\textbf{Answer.} Yes, in two ways, and the weights are themselves an example of the principle.

During learning the weights are adjusted in many small steps, each according to a small batch of
examples and therefore noisy. What occurs often pushes the weights again and again in the same
direction and prevails; the noise averages out as $1/\sqrt N$. What occurs rarely is lost in the noise
of the steps. In addition there are deliberately built-in brakes: weights are kept small, and training
is stopped in time. A noise limit is thereby set on purpose: fine structure that is not attested often
enough should not be learned, because it is mostly noise.

In storage each individual weight is surprisingly imprecise. With 16 bits this amounts to about $2.4$
to $3.3$ decimal digits \KM, and weights are often reduced afterwards to 8 or 4 bits without the model
becoming noticeably worse. This works only because knowledge does not reside in individual weights but
is distributed over very many at once. The network is thus itself a redundant, error-tolerant code,
like language. Good training solutions have exactly this property: a small shift of the weights
hardly changes the result \QM\,\cite{Hochreiter1997}. The flip side is the same as for language: exact
calculations that require full precision at every position succeed poorly inside the network.

% ============================================================
\section{Does anything work at all without set framework conditions?}
% ============================================================

\textbf{Question.} Does this mean that without set framework conditions nothing works anyway?

\textbf{Answer.} Essentially yes; for learning from data this is proved. According to the
no-free-lunch theorem of learning theory, averaged over all conceivable tasks no learning procedure is
better than any other, not even better than guessing \QM\,\cite{Wolpert1996}. A procedure is better
only if its built-in assumptions fit the actual world. Every procedure has such assumptions, even if
they are not stated: Fourier analysis presupposes periodicity, an image network that neighbouring
pixels belong together, a language model that text is a sequence with structure. Doc.~359
Theorem~D shows the difference in numbers: without a framework about 84\,000 examples, with the
appropriate framework at most 407, and from positive examples alone 30 \KM. Without a prior every
pattern is as good as any other.

Two qualifications belong to this. First, a framework can be weak when data are plentiful: the more
data, the less needs to be specified. Language models have a rather general framework and compensate
for it with very large amounts of text. The sparser the data, the stronger and more precise the
framework must be; with a few dozen constants of nature it must carry practically the whole structure.
This is the situation of FFGFT.

Second, a framework does not work merely because it has been set. A wrong or too broad framework also
adapts and then finds false patterns. The Balmer test of Section~\ref{sec:nulltest} shows both sides:
the narrow framework of fractions finds the true rule; the broad framework with 23 million expressions
fits everything and therefore no longer distinguishes anything \KM. A framework must therefore not
only be set but also tested; the null test measures whether it achieves more than chance.

For physics this means: every theory has settings. The Standard Model has its gauge groups and about
two dozen measured parameters; FFGFT has $T^4/\mathbb Z_3$, $\xi$ and the integer structural choices.
The question is not whether there are settings, but how few there are, how much follows from them,
and whether what follows withstands a test against chance: setting, derivation, comparison, null
test.

% ============================================================
\section{The common core in a single picture}
% ============================================================

All answers lead to the same structure. On the discrete side stand counting, words, group orders,
winding numbers and Born weights; there are gaps between the valid values, and noise smaller than such
a gap is absorbed. On the continuous side stand measured values and the weights of a network; there
every neighbouring value is valid.

\begin{center}
\begin{tikzpicture}[
  box/.style={draw, rounded corners, align=center, minimum width=5.2cm, minimum height=1.6cm, font=\small},
  arr/.style={-{Stealth[length=3mm]}, thick}]
\node[box] (D) at (0,0) {\textbf{Discrete}\\counting, words, group orders,\\winding numbers, Born weight $2/9$\\\emph{gap absorbs noise}};
\node[box] (K) at (8,0) {\textbf{Continuum}\\measured values, weights of a network\\[2pt]\emph{every neighbouring value valid}};
\draw[arr] (D.north east) to[bend left=18] node[above, font=\small, align=center]{measurement: structure is mapped\\\textbf{the noise limit sits here}} (K.north west);
\draw[arr] (K.south west) to[bend left=18] node[below, font=\small, align=center]{output of a network: a word is chosen\\noise below the gap is cut off} (D.south east);
\end{tikzpicture}
\end{center}

\noindent A network computes continuously inside and outputs discrete words; in the step from the
continuum to the discrete the noise is cut off as long as it is smaller than the gap between two
words. In measurement it runs the other way: an exact, discrete structure is mapped onto the
continuum, and there no gap absorbs the noise. The four core statements of this document follow:

\begin{enumerate}
\item Discrete structure is error-tolerant; the noise limit sits at the transition from the discrete
  to the continuum \BM.
\item Structure that is densely and variously attested in the data is learned; structure from a few
  exact values is not found from data, not even by a network \KM.
\item Such structure must be specified and then tested: setting, derivation, comparison with the
  measured value, null test. This is the path FFGFT takes \SM.
\item No learning without a framework, but a framework is only as good as its test \QM\,\cite{Wolpert1996}, \KM.
\end{enumerate}

\vspace{1em}
\noindent\textit{Verification scripts:} \texttt{2/python/Dok393\_Skripte/} ---
\texttt{pruef\_393a\_nulltest\_grammatik.py} (8/8, result log \texttt{pruef\_393a\_ergebnis.txt}),
\texttt{pruef\_393b\_zahlen.py} (11/11); together \texttt{run\_all\_393.sh}, 19/19.

\section*{Tool and method}
\addcontentsline{toc}{section}{Tool and method}

Claude (Anthropic) was used as a tool for: formulating the answers to the author's questions, writing
and running the verification scripts, and preparing the \LaTeX\ sources. The sections on how AI
assistants work describe the general operation of such systems; they do not rest on an inspection of
the model's own weights. Every number in the text is secured by a check in the scripts named; literature
references are marked \QM.

\begin{thebibliography}{99}
\bibitem{Balmer1885} J.~J.~Balmer, \textit{Notiz über die Spectrallinien des Wasserstoffs}, Annalen der Physik \textbf{261} (1885) 80--87.
\bibitem{Shannon1951} C.~E.~Shannon, \textit{Prediction and entropy of printed English}, Bell System Technical Journal \textbf{30} (1951) 50--64.
\bibitem{Hochreiter1997} S.~Hochreiter, J.~Schmidhuber, \textit{Flat minima}, Neural Computation \textbf{9} (1997) 1--42.
\bibitem{Rahaman2019} N.~Rahaman et al., \textit{On the spectral bias of neural networks}, Proceedings of the 36th International Conference on Machine Learning, PMLR \textbf{97} (2019) 5301--5310.
\bibitem{Udrescu2020} S.-M.~Udrescu, M.~Tegmark, \textit{AI Feynman: A physics-inspired method for symbolic regression}, Science Advances \textbf{6} (2020) eaay2631.
\bibitem{Wolpert1996} D.~H.~Wolpert, \textit{The lack of a priori distinctions between learning algorithms}, Neural Computation \textbf{8} (1996) 1341--1390.
\bibitem{ShalevShwartz2014} S.~Shalev-Shwartz, S.~Ben-David, \textit{Understanding Machine Learning: From Theory to Algorithms}, Cambridge University Press (2014).
\bibitem{NIST} A.~Kramida, Yu.~Ralchenko, J.~Reader, NIST ASD Team, \textit{NIST Atomic Spectra Database}, National Institute of Standards and Technology, Gaithersburg (hydrogen lines H$\alpha$--H$\delta$, air wavelengths).
\end{thebibliography}
